\chapter{Introdução}
\label{chap:introducao}

\section{Enquadramento e motivação}

A preparação de uma equipa de basquetebol exige decisões sobre o cinco
inicial, a distribuição de minutos e as rotações, conciliando o desempenho
imediato com a disponibilidade do plantel para jogos posteriores. Os dados
históricos, as estatísticas individuais e os registos \emph{play-by-play}
permitem analisar resultados e alinhamentos, apoiando a comparação de
alternativas pelo treinador.

Desenvolvido nas unidades curriculares de Aprendizagem Automática (AAUTIA)
e Planeamento em Inteligência Artificial (PLANTIA), este trabalho propõe
um sistema inteligente de apoio à decisão que articula previsões de
resultados e desempenho com a seleção de ações sujeitas a restrições.

\section{Problema de investigação}

Prever a probabilidade de vitória não basta para recomendar decisões:
é necessário representar os jogadores elegíveis, os alinhamentos, os
minutos de utilização e os objetivos do treinador. As previsões devem
ser compatíveis com as alternativas avaliadas pelo planeador, pois
estatísticas agregadas podem ser insuficientes para comparar alterações
concretas na equipa.

A questão central é: \emph{em que medida a integração de previsões de
aprendizagem automática com métodos de planeamento e otimização pode
melhorar as decisões simuladas de gestão de uma equipa de basquetebol,
face a estratégias fixas ou heurísticas simples?}

A avaliação distinguiu a qualidade preditiva da qualidade das decisões.
As associações históricas não demonstram efeitos causais de substituições
ou alterações de minutos; as conclusões dependeram das hipóteses e da
validade do ambiente de simulação.

\section{Objetivos do trabalho}

O objetivo geral foi desenvolver e avaliar um protótipo que integre previsão
e planeamento. Os objetivos específicos foram:

\begin{enumerate}
  \item Organizar dados históricos e construir indicadores com informação
  disponível no momento de cada previsão.
  \item Comparar modelos supervisionados para prever vitórias e, quando
  possível, o desempenho de alinhamentos.
  \item Analisar os fatores que contribuem para as previsões e facilitar
  a sua interpretação.
  \item Formular a seleção de jogadores e a distribuição de minutos sob
  restrições de disponibilidade, utilização e composição da equipa.
  \item Comparar métodos de otimização e planeamento num ambiente comum,
  incluindo sequências de vários jogos.
  \item Avaliar a integração dos módulos quanto à previsão, ao cumprimento
  das restrições e aos resultados simulados.
\end{enumerate}

\section{Abordagem proposta e âmbito}

A componente de AAUTIA prepara dados históricos da NBA e compara modelos
para prever vitórias, produção individual e margem final associada aos
cinco iniciais. A seleção dos modelos utiliza validação temporal, seguida
de avaliação numa temporada de teste. Os alinhamentos considerados são
os titulares, não todos os quintetos utilizados durante o jogo.

A componente de PLANTIA utiliza programação por restrições para construir
rotações, um MDP de horizonte finito para reorganizar blocos dentro de cada
período, comparando a solução exata com um algoritmo genético, e uma
seleção conjunta de rotações para vários jogos. As previsões
de vitória permitem quantificar objetivos competitivos, enquanto as
prioridades individuais alimentam funções de utilidade explicitamente
heurísticas.

O relatório descreve a implementação e as experiências registadas no
projeto. Distingue resultados preditivos, utilidade simulada e cumprimento
de restrições, sem interpretar estes últimos como provas de ganhos em
jogos reais. O âmbito é retrospetivo, com hipóteses de disponibilidade,
contexto e fadiga documentadas.

\section{Organização do relatório}

O Capítulo~\ref{chap:estado-arte} apresenta os trabalhos que fundamentam
a proposta e o Capítulo~\ref{chap:data} descreve a componente de dados.
O Capítulo~\ref{chap:metodologia} estabelece o protocolo e as interfaces.
Os Capítulos~\ref{chap:previsao-resultados} e~\ref{chap:desempenho}
desenvolvem as tarefas dos módulos de previsão de resultados e de avaliação de jogadores e alinhamentos. Os
Capítulos~\ref{chap:rotacoes}, \ref{chap:substituicoes} e~\ref{chap:horizonte}
apresentam o planeamento dos módulos de rotações, substituições e horizonte de vários jogos. O
Capítulo~\ref{chap:experiencias} reúne os resultados e a discussão;
o Capítulo~\ref{chap:conclusoes} apresenta conclusões e trabalho futuro.
